Few methods currently support online spike sorting for high-channel-count
probes. To the best of our knowledge, there is currently no readily-usable
online spike sorting solution available for chronic Neuropixels recordings in
freely moving mice.

We will implement a new online and adaptive spike-sorting method using the
hybrid strategy described in Section~Bonsai and Bonsai.ML.
%
Periodically, as new chunks of recording become available, the slower process
in this strategy will update a catalogue of waveforms of detected units, as
indicated in Section~Offline Spike Sorting.
%
With this catalogue, the lightweight process will perform template matching
(i.e., detect spikes and match the corresponding waveform to that of a unit in
the catalogue).
%
It will implemented as a Bonsai observable, and will operate on short chunks
of tens to hundreds of milliseconds.
%
Since several template matching algorithms are available, we will benchmark
them in terms of speed and accuracy against the offline spike sorting results.

As online spike sorting will become essential for future adaptive experiments,
we will integrate our solution into SWC and AIND Bonsai workflows for
real-time visualization and spike-based analysis.

\subsubsubsection{Deliverables}

\begin{enumerate}

    \item A new online spike-sorting method robust to non-stationarities in
        NaLoDuCo experimentation, targeting a latency in the order of tens to
        hundreds of milliseconds.

    \item Integration of the new online spike sorting method into existing
        Bonsai workflows at the SWC and AIND for closed-loop control and
        real-time experimental feedback.

    \item Open-source dissemination of code and tools for real-time spike
        sorting, including documentation and example Bonsai workflows using
        NaLoDuCo data.

\end{enumerate}

